
For the function $ f$ to be a density function of a continuous normal distribution we need the
following to hold

\[
    f_{\mu, \Sigma} (x,y) = \frac{1}{\sqrt{(2\pi)^2 \det \Sigma}} e^{- \frac{1}{2} ((x,y) - \mu)^{T}
    \Sigma^{-1} ((x,y)-\mu)} 
.\] 

We see that in the function $ h$  given to us there are no terms of order 1 meaning that it is most likely
that the expectation value vector $ \mu$  is zero. Through trial and error, we find that 

\[
    ((x,y))^{T} \Sigma^{-1} ((x,y)) = 2(x^2 + y^2 + xy)
\] 
where 
\[
    \Sigma^{-1} = \begin{bmatrix}
	2 & 1\\
	1 & 2
    \end{bmatrix}
.\] 
So our guess that the expectation value vector $ \mu $  is zero is correct since this matrix $ \Sigma$ 

\[
    \Sigma = \frac{1}{\det \Sigma^{-1} } \begin{bmatrix}
	2 & -1 \\
	-1 & 2
    \end{bmatrix} = \frac{1}{3} \begin{bmatrix}
	2 & -1 \\
	-1 & 2
    \end{bmatrix}
\] 
is symmetric and strictly positive definite so applying Theorem 14.9 (Density function of a
multidimensional normal distribution) we get that 

\[
    f(x,y) = f_{0, \Sigma} (x,y) = \frac{1}{\sqrt{(2\pi)^{2} \det(\Sigma)}} e^{- \frac{1}{2} ((x,y)-0)^{T} \Sigma^{-1}
    ((x,y)-0)} 
\] 
where we let $ C = 1/2\pi$.
In resume, we have shown that the function $ f$ is a density function of a random vector 
$ (X,Y)$ whose distribution is the 2-dimensional normal distribution with $ \mu(X,Y) = 0$ and $ \mathrm{Cov} = \Sigma$.
